\documentclass[12pt]{article}
\usepackage[utf8]{inputenc}
\usepackage[T1]{fontenc}
\usepackage[margin=1in]{geometry}
\usepackage{setspace}
\usepackage{xcolor}
\usepackage[colorlinks=true, urlcolor=blue!60!black, linkcolor=blue!60!black]{hyperref}
\setlength{\parindent}{0em}
\setlength{\parskip}{0.6em}
\renewcommand{\baselinestretch}{1.0}
\begin{document}
\begin{center}
{\LARGE \textit{Louisiana v. Callais}, 608 U.S. \_\_\_ (2026)\par}
\vspace{0.5em}
{\large Case Brief --- Two-Page Edition\par}
\end{center}

\section{Memory Jogger}
Louisiana's second majority-Black district was a racial gerrymander: Section~2, now requiring a strong inference of intentional discrimination, did not require it (\href{https://supreme.justia.com/cases/federal/us/608/24-109/}{Justia}).

\section{Detailed Case Facts}
After a court found Louisiana's one-majority-Black-district map likely violated Section~2, the Legislature enacted SB8, adding District~6, a roughly 250-mile diagonal from Shreveport to Baton Rouge drawn to protect Speaker Johnson, Majority Leader Scalise, and Representative Letlow. ``Non-African American'' voters led by Phillip Callais sued, calling it a racial gerrymander.

\section{Procedural History}
A three-judge court enjoined SB8, 2--1. The Supreme Court stayed the injunction, heard argument, ordered reargument on the constitutional question, and affirmed on April 29, 2026.

\section{Judicial Votes}
Affirmed, 6--3. Alito wrote for Roberts, Thomas, Gorsuch, Kavanaugh, and Barrett; Thomas (with Gorsuch) concurred; Kagan (with Sotomayor and Jackson) dissented.

\section{Holding}
Section~2 compliance can be a compelling interest, but Section~2 imposes liability only on a strong inference of intentional discrimination. Plaintiffs must offer a race-blind map meeting all the State's goals, including partisan ones, show polarization controlling for party, and rely on current evidence. The \textit{Robinson} plaintiffs failed, so SB8 fell.

\section{Analysis of Opinions}
\textbf{Majority:} Because the Fifteenth Amendment bars only purposeful discrimination, \textit{Boerne} limits Section~2 to intent; Alito ``updated'' \textit{Gingles} for social progress, a two-party South, \textit{Rucho}, and computing power, distinguishing \textit{Allen v. Milligan}. \textbf{Thomas:} Section~2 should not reach districting at all. \textbf{Kagan:} the ruling restores the \textit{Bolden} intent test Congress rejected in 1982, flouts stare decisis, and immunizes textbook dilution.

\section{Examples: Future Applications}
\textbf{Same-Side:} a legislature cracks a Black urban district to elect an all-Republican delegation; the claim fails. \textbf{Opposite-Side:} a one-party city splits a Hispanic neighborhood despite race-blind maps meeting every criterion. \textbf{Fence-Sitter:} emails show racial hostility, but a partisan goal requires dismantling the district.

\section{Critique}
Brian Fitzpatrick defends the decision as color-blind equal protection, though even sympathizers call it awkward and evasive about \textit{Milligan}. Erwin Chemerinsky and the Brennan Center call it the judicial nullification of Section~2, letting States immunize dilution by declaring a partisan gerrymander.

\section{Key Quotations}
\begin{quote}
``In short, \S~2 imposes liability only when the evidence supports a strong inference that the State intentionally drew its districts to afford minority voters less opportunity because of their race. Not only does this interpretation follow from the plain text of \S~2, but it is consistent with the limited authority that the Fifteenth Amendment confers.''
\end{quote}
\noindent --- Justice Alito, Opinion of the Court (slip op., at 26).

\end{document}
